% Declare any packages and options you are using here.

% In the ACM template, there is a list of allowed packages and list of packages
% already included as part of acmart. See here:
% https://authors.acm.org/proceedings/production-information/accepted-latex-packages

% INCLUDED PACKAGES %%%%%%%%%%%%%%%%%%%%%%%%%%%%%%%%%%%%%%%%%%%%%%%%%%%%%%%%%%%%
% The acmart template includes some packages already. DO NOT UNCOMMENT ANYTHING
% IN THESE LISTS - they're just there to help you remember which ones you're
% using that are included with acmart.
% Enable features useful for typesetting mathematical notation (e.g., multiline
% alignment for equations).
%\usepackage{amsmath}
% Balance the final page columns.
%\usepackage{balance}
% Make tables look nicer.
%\usepackage{booktabs}
% Allow figures
%\usepackage{graphicx}
% Create hypertext links for citations and references.
%\usepackage{hyperref}
% Define colors for comments, etc.
%\usepackage{xcolor}
%%%%%%%%%%%%%%%%%%%%%%%%%%%%%%%%%%%%%%%%%%%%%%%%%%%%%%%%%%%%%%%%%%%%%%%%%%%%%%%%

% ACTIVE PACKAGES %%%%%%%%%%%%%%%%%%%%%%%%%%%%%%%%%%%%%%%%%%%%%%%%%%%%%%%%%%%%%%
% Used packages. Keep packages in alphabetical order, except for hyperref below.
%
% Use acronyms (define actual acronyms in setup/acro.tex). The printonlyused
% option is probably not useful in papers since an acronym list isn't printed at
% all, but leave it here just in case.
\usepackage[printonlyused]{acronym}
% Highly customizable enumeration, including inline lists.
\usepackage[inline]{enumitem}
% For spacing after user-defined commands and others.
\usepackage{xspace}
%%%%%%%%%%%%%%%%%%%%%%%%%%%%%%%%%%%%%%%%%%%%%%%%%%%%%%%%%%%%%%%%%%%%%%%%%%%%%%%%

% OTHER PACKAGES %%%%%%%%%%%%%%%%%%%%%%%%%%%%%%%%%%%%%%%%%%%%%%%%%%%%%%%%%%%%%%%
% Packages that are not used, but may be useful in certain situations. Move and
% uncomment to the list above to use.
%
% Use acronyms (define actual acronyms in setup/acro.tex).
%\usepackage{acronym}
%
% Enable algorithm descriptions in pseudocode.
%\usepackage{algorithmic}
%
% Enable features useful for typesetting mathematical notation (e.g., multiline
% alignment for equations).
%\usepackage{amsmath}
%
% Load AMS fonts.
%\usepackage{amsfonts}
%
% Enable commands for certain mathematical symbols (e.g., \varnothing for the
% empty set).
%\usepackage{amssymb}
%
% Get bitmap size and resolutions
%\usepackage{bmpsize}
%
% Make tables look nicer
%\usepackage{booktabs}
%
% Compress and sort multi-citations. You _shouldn't_ use this in most modern
% templates since they have their own style, but it can come in handy.
%\usepackage{cite}
%
% Highly customizable enumeration, including inline lists.
%\usepackage[inline]{enumitem}
%
% Enable message sequence charts, which are useful for showing messages sent in
% a particular order (e.g., in network protocols).
%\usepackage{msc}
%
% Enable multiple subfigures within a single figure.
%\usepackage{subfig}
%
% Enable table columns to fill/wrap a specific width. This isn't allowed by ACM
% templates, so your source might get rejected by TAPS if you use it.
%\usepackage{tabularx}
%
% Enable various text symbols (e.g., currency symbols, bullet points, copyright
% symbols, musical notation).
%\usepackage{textcomp}
%
% Make URLs display correctly in citations. The hyphens option allows line
% breaks after hyphens.
%\usepackage[hyphens]{url}
%
% For spacing after user-defined commands and others.
%\usepackage{xspace}
%%%%%%%%%%%%%%%%%%%%%%%%%%%%%%%%%%%%%%%%%%%%%%%%%%%%%%%%%%%%%%%%%%%%%%%%%%%%%%%%
